\documentclass[12pt]{article}
\usepackage[utf8]{inputenc}
\usepackage[T1]{fontenc}
\usepackage{hyperref}
\title{Maasina Fulfulde Language Resources for K-12 Curriculum Development}
\author{Compiled by Claude Sonnet 5.5 (claude-sonnet-5-5) for Andr\'as Kornai}
\date{}
\begin{document}
\maketitle

\section{Introduction}
Maasina Fulfulde (ISO 639-3 ffm) is spoken in Mali. Its Latin alphabet was standardised in 1967 (32 letters, including implosives and the glottal stop); it was historically written in Arabic script. The searches found no Malian school primers or ministry textbooks, only general reference and learner materials.

\section{Primary Educational Resources}
No school primers or textbooks were found in the search results.

\section{Linguistic and Reference Materials}
\begin{itemize}
\item Maasina Fulfulde (Wikipedia): \url{https://en.wikipedia.org/wiki/Maasina_Fulfulde}
\item Fula alphabets (Wikipedia): \url{https://en.wikipedia.org/wiki/Fula_alphabets}
\item Fulfulde Maasina--English--French lexicon (Michigan State University Press), for learners at intermediate or advanced level: \url{https://msupress.org/9780870139420/fulfulde-maasina-english-french-lexicon}
\end{itemize}

\section{Religious and Community Resources}
None verified in the search results.

\section{Digital and Audio Resources}
\begin{itemize}
\item Fulfulde language learning resources (University of Wisconsin LCTL resources), including the Peace Corps Mali Fulfulde technical manual: \url{https://wisc.pb.unizin.org/lctlresources/chapter/fulfulde-language-learning-resources/}
\item Fula Language on the Internet (Columbia University Libraries), which points to further materials: \url{https://library.columbia.edu/libraries/global/virtual-libraries/african_studies/languages/fula.html}
\end{itemize}

\section{Curriculum Development Context}
Although the language is used in education in Mali, no official curriculum documents were located in these searches. Follow-up searches of the Malian Ministry of Education and SIL Mali publications are recommended.

\end{document}